\documentclass[journal]{vgtc}                % final (journal style)
%\documentclass[review,journal]{vgtc}         % review (journal style)
%\documentclass[widereview]{vgtc}             % wide-spaced review
%\documentclass[preprint,journal]{vgtc}       % preprint (journal style)

%% Uncomment one of the lines above depending on where your paper is
%% in the conference process. ``review'' and ``widereview'' are for review
%% submission, ``preprint'' is for pre-publication, and the final version
%% doesn't use a specific qualifier.

%% Please use one of the ``review'' options in combination with the
%% assigned online id (see below) ONLY if your paper uses a double blind
%% review process. Some conferences, like IEEE Vis and InfoVis, have NOT
%% in the past.

%% Please note that the use of figures other than the optional teaser is not permitted on the first page
%% of the journal version.  Figures should begin on the second page and be
%% in CMYK or Grey scale format, otherwise, colour shifting may occur
%% during the printing process.  Papers submitted with figures other than the optional teaser on the
%% first page will be refused. Also, the teaser figure should only have the
%% width of the abstract as the template enforces it.

%% These few lines make a distinction between latex and pdflatex calls and they
%% bring in essential packages for graphics and font handling.
%% Note that due to the \DeclareGraphicsExtensions{} call it is no longer necessary
%% to provide the the path and extension of a graphics file:
%% \includegraphics{diamondrule} is completely sufficient.
%%
\ifpdf%                                % if we use pdflatex
  \pdfoutput=1\relax                   % create PDFs from pdfLaTeX
  \pdfcompresslevel=9                  % PDF Compression
  \pdfoptionpdfminorversion=7         % create PDF 1.7
  \ExecuteOptions{pdftex}
  \usepackage{graphicx}  
  \graphicspath{ {pictures/} }              % allow us to embed graphics files
  \DeclareGraphicsExtensions{.pdf,.png,.jpg,.jpeg} % for pdflatex we expect .pdf, .png, or .jpg files
\else%                                 % else we use pure latex
  \ExecuteOptions{dvips}
  \usepackage{graphicx}                % allow us to embed graphics files
  \DeclareGraphicsExtensions{.eps}     % for pure latex we expect eps files
\fi%

%% it is recomended to use ``\autoref{sec:bla}'' instead of ``Fig.~\ref{sec:bla}''
\graphicspath{{figures/}{pictures/}{images/}{./}} % where to search for the images

\usepackage{microtype}                 % use micro-typography (slightly more compact, better to read)
\PassOptionsToPackage{warn}{textcomp}  % to address font issues with \textrightarrow
\usepackage{textcomp}                  % use better special symbols
\usepackage{mathptmx}                  % use matching math font
\usepackage{times}                     % we use Times as the main font
\renewcommand*\ttdefault{txtt}         % a nicer typewriter font
\usepackage{cite}                      % needed to automatically sort the references
\usepackage{tabu}                      % only used for the table example
\usepackage{booktabs}                  % only used for the table example
\usepackage[section]{placeins}
\usepackage{float}
\usepackage{hyperref}
\usepackage[all]{hypcap}
\usepackage{enumitem} 
\hypersetup{
    colorlinks = true,
    citecolor = {black},
    urlcolor = {red}
}

%% We encourage the use of mathptmx for consistent usage of times font
%% throughout the proceedings. However, if you encounter conflicts
%% with other math-related packages, you may want to disable it.

%% In preprint mode you may define your own headline.
%\preprinttext{To appear in IEEE Transactions on Visualization and Computer Graphics.}

%% If you are submitting a paper to a conference for review with a double
%% blind reviewing process, please replace the value ``0'' below with your
%% OnlineID. Otherwise, you may safely leave it at ``0''.
\onlineid{0}

%% declare the category of your paper, only shown in review mode
\vgtccategory{Research}
%% please declare the paper type of your paper to help reviewers, only shown in review mode
%% choices:
%% * algorithm/technique
%% * application/design study
%% * evaluation
%% * system
%% * theory/model

%% Paper title.
\title{An Evaluation of Poemage}

%% This is how authors are specified in the journal style

%% indicate IEEE Member or Student Member in form indicated below
\author{Daniel Schroeder, Will Sims, and Bungo Takahashi}
\authorfooter{
%% insert punctuation at end of each item
\item
 Daniel Schroeder is a computer science student at Oregon State University. E-mail: schrodan@oregonstate.edu.
\item
 BUNGO ENTER HERE
\item
 WILL ENTER HERE
}

%other entries to be set up for journal
%\shortauthortitle{Biv \MakeLowercase{\textit{et al.}}: Crimes on Oregon College Campuses}
%\shortauthortitle{Firstauthor \MakeLowercase{\textit{et al.}}: Paper Title}

%% Abstract section.
\abstract{This paper is a visualization evaluation of the program Poemage, created by Nina McCurdy, Julie Lein, Katharine Coles, and Miriah Meyer from the University of Utah. We will evaluate the effectiveness of this visualization as well as suggest future implementations that could enhance the application. %
} % end of abstract

%% Keywords that describe your work. Will show as 'Index Terms' in journal
%% please capitalize first letter and insert punctuation after last keyword
\keywords{Visualization in the humanities, design studies, text and document data, graph/network data}

%% ACM Computing Classification System (CCS). 
%% See <http://www.acm.org/class/1998/> for details.
%% The ``\CCScat'' command takes four arguments.

\CCScatlist{ % not used in journal version
 \CCScat{K.6.1}{Management of Computing and Information Systems}%
{Project and People Management}{Life Cycle};
 \CCScat{K.7.m}{The Computing Profession}{Miscellaneous}{Ethics}
}

%% Uncomment below to include a teaser figure.
%%\teaser{
%%  \centering
%%  \includegraphics[width=\linewidth, height=3.5in]{FF}
%%  \caption{Photo by: Zach Freedel. The photo of the fog was taken in Corvallis, Oregon, on the campus of Oregon State University.}
%%	\label{fig:teaser}
%%}

%% Uncomment below to disable the manuscript note
\renewcommand{\manuscriptnotetxt}{}

%% Copyright space is enabled by default as required by guidelines.
%% It is disabled by the 'review' option or via the following command:
% \nocopyrightspace

\vgtcinsertpkg

%%%%%%%%%%%%%%%%%%%%%%%%%%%%%%%%%%%%%%%%%%%%%%%%%%%%%%%%%%%%%%%%
%%%%%%%%%%%%%%%%%%%%%% START OF THE PAPER %%%%%%%%%%%%%%%%%%%%%%
%%%%%%%%%%%%%%%%%%%%%%%%%%%%%%%%%%%%%%%%%%%%%%%%%%%%%%%%%%%%%%%%%

\begin{document}

%% The ``\maketitle'' command must be the first command after the
%% ``\begin{document}'' command. It prepares and prints the title block.

%% the only exception to this rule is the \firstsection command
\firstsection{Introduction}

\maketitle

The work presented here yields the following benefits and contributions:
\begin{itemize}
\setlength\itemsep{0em}
\item blah 
\item blah
\item blah
\end{itemize}

\section{Related Work} \label{research}

Close reading through computational means is a very difficult problem to tackle due to its large problem set and complexity by nature. Most design decisions for computational close reading tend to focus on a specific area of interest within close reading, for example sound or rhyme. Nina McCurdy from The University of Utah wrote: ``Close reading covers a broad range of tasks, encompasses varying styles of analysis, allows many different points of entry, and accepts an extensive range of sometimes radically divergent interpretations'' when describing the complexity of text analysis \cite{poemage}.\\
\indent There have been numerous previous projects trying to tackle close reading via visualization and computation like E-Margin \cite{emargin}, Serendip \cite{serendip}, and Myopia \cite{myopia}. The Poemage team took a slightly different course during their development by focusing on sounds and creating a program that could ``automatically sonify a poem'' \cite{poemage}. All previous projects address the problem of close reading as an interpretation of text at the level of individual words in order to interpret and examine literary context. A large issue with addressing this problem is that it is usually done in an academic setting where someone is trying to ``teach'' the literature, and hand written evaluations tend to be messy and difficult to understand. These close reading visualization projects attempt to help visualize textual interpretations so that it can be taught and evaluated easily and effectively.\\
\indent We wanted to branch off from the Poemage style of using drawn ellipses to highlight words of interest and search for other methods of text highlighting that may increase the effectiveness of the visualization. We sought to highlight individual syllables or letters that applied to the filtered rhyme scheme or rhyme rule rather than circling the entire word. There has been extensive research on text highlighting techniques conducted by the information visualization community. A paper from IEEE Transactions on Visualization and Computer Graphics states that ``Text highlighting is important in any scenario where close reading (sequential word-by-word reading) is required and text annotations exist, that should be made accessible to the reader'' \cite{texthighlighting}. We wanted to explore the possibilities of using background coloring to target different parts of a single word to beautify the user interface as a whole and enhance the effectiveness of the Poemage program. By using different background colors used to target specific letters and syllables, the user will be able to clearly identify the sonic devices being targeted and be able to see where in the word it is present. The issue of where in the word the sonic device is located is not expressed through the use of drawn ellipses in the current version of Poemage, and users become overpowered with clutter and complexity when multiple filters are applied. 



\section{Method} \label{method}
\subsection{Method Overview}

\subsection{Usability Evaluation}

\subsection{Poemage With Different Languages}

\section{Enhancements} \label{enhancements}

\section{Results and Performance} \label{results}
\subsection{Method Overview}

\subsection{Usability Evaluation}

\subsection{Poemage With Different Languages}

\subsection{Performance Analysis}

\section{Conclusions and Future Work} \label{conclusion}

\acknowledgments{
The authors wish to thank Nina McCurdy and the faculty at The University of Utah that built the Poemage program as well as professor Eugene Zhang and teaching assistant Islam Al Musaly.}

\newpage
%\bibliographystyle{abbrv}
\bibliographystyle{abbrv-doi}
%\bibliographystyle{abbrv-doi-narrow}
%\bibliographystyle{abbrv-doi-hyperref}
%\bibliographystyle{abbrv-doi-hyperref-narrow}

\bibliography{refs}
\end{document}

